\section{近端参数的选择}\label{sec:parameters}

前文先固定 $\lambda$，再检验次单调性与正则性的补偿。
Luke--Tam~\cite{LukeTam2025} 也将缩放算子 $\lambda F$ 的性质放在假设中。
本节反过来，从原算子的信息确定可用参数，最后允许每步更换 $\lambda_k$。
改变参数会改变近端对，故下列共同估计均须对所考虑的全部参数有效。

\subsection{DD1 与 DD2 的参数反演}

先设 $\ell>0$、$\sigma\ge0$ 与参数无关，且每个所考虑的近端对满足
\begin{equation}\label{eq:parameters-common-linear}
 d(y)\le\ell r_F(y),\qquad
 \exists p\in P_S(x):\quad\ip{x-y}{y-p}\ge-\sigma d(x)^2.
\end{equation}

\begin{proposition}[DD1 的参数条件]\label{prop:parameters-dd1}
给定 $0<r<1$，令
\begin{equation}\label{eq:parameters-linear-target}
 \lambda_r=\frac\ell r
 \begin{cases}
 \sqrt{1-r^2+2\sigma},&\sigma\le r(1+r),\\[2pt]
 (1+\sqrt{1+4\sigma})/2,&\sigma>r(1+r).
 \end{cases}
\end{equation}
则 DD1 中保证 $d(y)\le r d(x)$ 的分段条件可写为 $\lambda\ge\lambda_r$。
存在某个统一收缩比例 $r<1$ 的 DD1 条件为
\begin{equation}\label{eq:parameters-linear-strict}
 \lambda>\ell
 \begin{cases}
 \sqrt{2\sigma},&0\le\sigma\le2,\\[2pt]
 (1+\sqrt{1+4\sigma})/2,&\sigma\ge2.
 \end{cases}
\end{equation}
\end{proposition}

\begin{proof}
令 $b_\sigma=(1+\sqrt{1+4\sigma})/2$。
将 $L=\ell/\lambda$ 代入 DD1 的两项收缩估计，条件等价于
\[
 \lambda\ge\frac{\ell b_\sigma}r
 \quad\text{或}\quad
 \lambda\ge\frac{\ell\sqrt{1+2\sigma-r^2}}r.
\]
取两个下界中较小者；两者在 $\sigma=r(1+r)$ 处相等，即得
\eqref{eq:parameters-linear-target}。令 $r\uparrow1$ 即得严格条件
\eqref{eq:parameters-linear-strict}。
\end{proof}

一般函数情形中，设同一 $F,S$ 对所有候选参数 $\lambda\in\Lambda$
及相应允许对 $y\in\mathcal J_\lambda(x)\subseteq J_\lambda(x)$ 均满足
\[
 \rho=r_F(y)\in D_\psi,\qquad d(y)\le\psi(\rho),\qquad
 \exists p\in P_S(x):\quad\ip{x-y}{y-p}\ge-D(t),\quad t=d(x),
\]
其中 $D\ge0$，$0\in D_\psi$。给定目标 $0<\varphi(t)<t$，记
\begin{equation}\label{eq:parameters-bad-layer}
 E_t=\{\rho\in D_\psi:\psi(\rho)>\varphi(t)\},\qquad
 q_t=\inf E_t,\qquad B_D(t)=\frac{t+\sqrt{t^2+4D(t)}}2.
\end{equation}

\begin{theorem}[DD2 的参数条件]\label{thm:parameters-dd2}
当 $E_t\ne\varnothing$ 且 $0<q_t<\infty$ 时，令
\begin{equation}\label{eq:parameters-general-threshold}
 H(t)=
 \begin{cases}
 \dfrac{\sqrt{t^2-\varphi(t)^2+2D(t)}}{q_t},
 &D(t)\le\varphi(t)(t+\varphi(t)),\\[6pt]
 \dfrac{B_D(t)}{q_t},&D(t)>\varphi(t)(t+\varphi(t)).
 \end{cases}
\end{equation}
DD2 保证 $d(y)\le\varphi(t)$ 的分段条件等价于 $\lambda\ge H(t)$；
仅当第二支且 $q_t\in E_t$ 时，应改为 $\lambda>H(t)$。
若 $E_t=\varnothing$，记 $H(t)=0$，不需附加参数条件；
若 $E_t\ne\varnothing$ 且 $q_t=0$，记 $H(t)=+\infty$，无有限正参数满足该判据。
因此若 $H_R:=\sup_{0<t\le R}H(t)<\infty$，则每个
$\lambda\in\Lambda$、$\lambda>H_R$ 均满足全部距离 $0<t\le R$ 的下降条件。
在 $\lambda=H_R>0$ 时须逐点检查上述端点条件。
\end{theorem}

\begin{proof}
在 DD2 中代入 $r=\varphi(t)/t$、$z=\lambda q_t/t$、$\delta=D(t)/t^2\ge0$。
第三支 $r(r-1)<0$ 不适用。中间支要求
$z\ge\sqrt{1-r^2+2\delta}$；该根不小于 $1-r$，且不大于 $1+r$
当且仅当 $\delta\le r(1+r)$，得到第一支。
若 $\delta>r(1+r)$，则由 $z(z-1)\ge\delta$ 解出第二支。
两支交点属于中间支，允许等号；其余开闭端点直接沿用 DD2 的规定。
空集情形由残差估计已得 $d(y)\le\varphi(t)$；非空且 $q_t=0$ 时 $z=0$，
分段式不可能容纳 $\delta\ge0$。最后对全部 $t$ 同时应用即可。
\end{proof}

若 $\psi$ 在从零开始的区间上连续严格递增，且目标值属于该限制的像的内部，
则 $q_t=\psi^{-1}(\varphi(t))$、$q_t\notin E_t$。
若 $\lambda>\limsup_{t\downarrow0}H(t)$，则在足够小的距离区间内满足该条件。
这些结论还须结合近端存在性、零输入驻定及相应长度和留域条件；
它们没有断言固定算子只有满足此判据才能收敛。

\begin{example}[小参数下的距离扩张]\label{ex:parameters-rotation}
令 $Q$ 为平面直角旋转，$F=-I+2Q$，$S=\{0\}$。则
\[
 I+\lambda F=(1-\lambda)I+2\lambda Q,\qquad
 \norm{(I+\lambda F)y}^2=(1-2\lambda+5\lambda^2)\norm y^2.
\]
因此每个 $\lambda>0$ 的近端步均唯一存在，而实际距离比为
\[
 \frac{d(y)}{d(x)}=\frac1{\sqrt{1-2\lambda+5\lambda^2}}.
\]
当 $0<\lambda<2/5$ 时距离增大，当 $\lambda>2/5$ 时距离缩小。
这里 $\ell=1/\sqrt5$，最小二次缺陷系数
$\sigma_\lambda=\lambda/(1-2\lambda+5\lambda^2)$ 虽趋于零，
仍不足以保证小参数下的收敛。
\end{example}

\subsection{原算子条件与变参数迭代}

以下条件在选取参数之前给定，同时保证共同邻域内的近端存在性与收敛。

\begin{theorem}[共同参数范围与变参数收敛]\label{thm:parameters-variable}
设 $F:\R^n\to\R^n$ 在 $V=\overline B(p_0,R_0)$ 的邻域连续，
$S\subseteq\zer F$ 非空闭，$p_0\in S$。设 $\ell>0$、$0\le\mu<1$，
对全部 $y\in V$、$p\in S\cap V$ 有
\begin{equation}\label{eq:parameters-raw-angle}
 d(y)\le\ell\norm{F(y)},\qquad
 \ip{F(y)}{y-p}\ge-\mu\norm{F(y)}\norm{y-p}.
\end{equation}
取 $0<r_0<\min\{R_0/2,\sqrt{1-\mu^2}R_0\}$，令
$U=B(p_0,r_0)$、$\mathcal J_\lambda(x)=J_\lambda(x)\cap B(p_0,R_0)$。
则每个 $\lambda>0$ 下，$\mathcal J_\lambda$ 在 $U$ 上非空。
给定 $0<r<1$，令
\begin{equation}\label{eq:parameters-raw-threshold}
 \lambda_*(r)=\ell
 \begin{cases}
 \mu+\sqrt{\mu^2+r^{-2}-1},&\mu\le r/\sqrt{1+r^2},\\[3pt]
 1/(r\sqrt{1-\mu^2}),&\mu\ge r/\sqrt{1+r^2}.
 \end{cases}
\end{equation}
当 $\lambda\ge\lambda_*(r)$ 时，每个允许对均满足
\begin{equation}\label{eq:parameters-common-estimates}
 d(y)\le r d(x),\qquad \norm{x-y}\le\frac{d(x)}{\sqrt{1-\mu^2}}.
\end{equation}
若 $x^0\in U$ 且
\begin{equation}\label{eq:parameters-local-start}
 \frac{d(x^0)}{(1-r)\sqrt{1-\mu^2}}<\dist(x^0,U^c),
\end{equation}
则任取 $\lambda_k\ge\lambda_*(r)$ 和
$x^{k+1}\in\mathcal J_{\lambda_k}(x^k)$，迭代可无限续接、始终留在 $U$，
并收敛于某个 $x^\infty\in S\cap U$，且
\begin{equation}\label{eq:parameters-variable-tail}
 d(x^k)\le r^k d(x^0),\qquad
 \sum_{j=k}^{\infty}\norm{x^{j+1}-x^j}
 \le\frac{d(x^k)}{(1-r)\sqrt{1-\mu^2}}.
\end{equation}
若连续性及~\eqref{eq:parameters-raw-angle} 全局成立，则可取完整关系
$\mathcal J_\lambda=J_\lambda$，上述结论对任意初值成立，无需~\eqref{eq:parameters-local-start}。
\end{theorem}

\begin{proof}
先证存在性。对 $H_\tau(y)=y+\tau\lambda F(y)$，在 $\partial V$ 上用 $p=p_0$
并记 $h=\tau\lambda\norm{F(y)}$，有
\[
 \norm{H_\tau(y)-p_0}^2\ge R_0^2+h^2-2\mu R_0h
 \ge(1-\mu^2)R_0^2.
\]
故对 $\norm{x-p_0}<\sqrt{1-\mu^2}R_0$，同伦在边界避开 $x$。
Brouwer 度的同伦不变性给出
$\deg(H_1-x,B(p_0,R_0),0)=\deg(I-x,B(p_0,R_0),0)=1$，
从而球内有根，且输入邻域与 $\lambda$ 无关。

对 $x\in U$，每个 $p_x\in P_S(x)$ 满足
$\norm{p_x-p_0}\le2\norm{x-p_0}<R_0$。
记 $b=\norm{y-p_x}$、$s=\norm{x-y}$、$t=d(x)$、$a=d(y)$，则
\begin{equation}\label{eq:parameters-angle-energy}
 t^2\ge b^2+s^2-2\mu bs
 =(b-\mu s)^2+(1-\mu^2)s^2.
\end{equation}
这已给出步长上界，并在 $t=0$ 时保证驻定。
当 $a>0$ 时，$h=s/a\ge\lambda/\ell$、$b/a\ge1$。
将右端对 $b/a\ge1$ 最小化，得
\[
 \frac{t^2}{a^2}\ge f(h):=
 \begin{cases}1+h^2-2\mu h,&\mu h\le1,\\
 (1-\mu^2)h^2,&\mu h\ge1.\end{cases}
\]
$f$ 在 $h\ge\mu$ 上非减，解 $f(h)=r^{-2}$ 即得
\eqref{eq:parameters-raw-threshold} 的两支；两式在交点相等。
故 $\lambda\ge\lambda_*(r)$ 给出 $a\le rt$；$a=0$ 时结论显然。

令 $C=(1-\mu^2)^{-1/2}$。每个输入仍在 $U$ 的有限迭代段均有
$t_{k+1}\le rt_k$、$s_k\le Ct_k$，因此连同最后输出在内，总长度不超过
$Ct_0/(1-r)$。由~\eqref{eq:parameters-local-start}，最后输出仍在 $U$。
逐步归纳即得无限续接与留域；从任意第 $k$ 步求和得
\eqref{eq:parameters-variable-tail}。有限长度保证柯西性，$t_k\to0$ 与 $S$ 闭
保证极限在 $S$，严格初值不等式保证极限在 $U$。
全局情形对任意输入选足够大的 $R_0$ 重复上述度论证明，
得到完整近端关系处处非空，再应用同样的估计即可。
\end{proof}

上述原算子条件也能直接接入 DD1。由~\eqref{eq:parameters-angle-energy}，
$t^2\ge2(1-\mu)bs$，从而
\[
 \ip{x-y}{y-p_x}\ge-\mu bs
 \ge-\frac{\mu}{2(1-\mu)}t^2.
\]
将共同系数 $\sigma=\mu/[2(1-\mu)]$ 代入
\eqref{eq:parameters-linear-target} 即可选参；保留角度信息则得到
\eqref{eq:parameters-raw-threshold}。这些结果允许逐步选择参数，
并未给出每个固定算子的最优参数。
